\documentclass[aspectratio=169,11pt]{beamer}

\usetheme{Madrid}
\usecolortheme{default}
\usefonttheme{professionalfonts}
\setbeamertemplate{navigation symbols}{}
\setbeamertemplate{footline}[frame number]

\usepackage{amsmath, amssymb, booktabs}

\title{Worker Adjustment, Skills, Training, And Career Transitions}
\subtitle{Supply-Side Adjustment Under Technological Change}
\author{Special Topic 5: Technology, Innovation, and Labor -- Week 3}
\date{}

\begin{document}

\begin{frame}
  \titlepage
\end{frame}

\begin{frame}{Central question}
\begin{itemize}
    \item When technology changes task prices and career risk, which workers can adjust?
    \item Which workers adapt through skills, training, mobility, or job redesign?
    \item Which workers bear persistent earnings, employment, and welfare losses?
    \item Week 3 discipline: aggregate reallocation is not the same object as worker welfare.
\end{itemize}
\end{frame}

\begin{frame}{From demand shocks to supply response}
\begin{itemize}
    \item Week 2: automation and AI shift labor demand across tasks, firms, occupations, and places.
    \item Week 3: exposed workers choose whether to learn, retrain, search, switch, move, bargain, or exit.
    \item The labor object is the cost and feasibility of adjustment.
    \item The policy object is not generic reskilling; it is incidence across workers with different constraints.
\end{itemize}
\end{frame}

\begin{frame}{Worker-adjustment framework}
\small
\[
V_{it}
=
w_{it}(s_{it}, o_{it}, f_{it}, \tau_t)
- c_i(a_{it})
- \kappa_i(m_{it})
+ E_t \sum_{r=t+1}^{T} \beta^{r-t} u_{ir}
\]
\begin{itemize}
    \item $s$: skills; $o$: occupation; $f$: employer; $\tau$: technology environment.
    \item $a$: costly adjustment effort such as training, coursework, search, or learning.
    \item $m$: mobility across employers, occupations, sectors, or places.
    \item Technology changes current pay, future skill values, and the option value of moving.
\end{itemize}
\end{frame}

\begin{frame}{Adjustment margins}
\small
\begin{tabular}{p{0.24\linewidth} p{0.36\linewidth} p{0.26\linewidth}}
\toprule
Margin & Labor object & Main question \\
\midrule
Early skill choice & major, field, vocational track & who enters resilient skill paths? \\
Within-firm adaptation & retraining, reassignment, learning & do incumbents adapt or get replaced? \\
Between-job mobility & employer, occupation, sector, location & can workers move to growing tasks? \\
Labor-force exit & unemployment, disability, retirement & is exit substituting for retraining? \\
\bottomrule
\end{tabular}
\vspace{0.35cm}
\begin{itemize}
    \item Staying can mean adaptation or being stuck; moving can mean opportunity or displacement.
\end{itemize}
\end{frame}

\begin{frame}{Aggregate adjustment vs worker losses}
\begin{itemize}
    \item A market can fill new vacancies while original exposed workers lose earnings for years.
    \item Firms can upgrade skill composition by hiring new workers instead of retraining incumbents.
    \item Average wages can rise through composition even when displaced workers are worse off.
    \item Local employment can recover through entry, migration, or retirement of exposed workers.
    \item The welfare question follows workers, not only jobs.
\end{itemize}
\end{frame}

\begin{frame}{Skill specificity and portability}
\begin{itemize}
    \item Portable skills retain value across firms, occupations, technologies, and local labor markets.
    \item Specific skills can pay well before a shock and depreciate sharply after tasks change.
    \item STEM and other technology-intensive careers can require repeated updating as tools change.
    \item Two equally exposed workers can face different losses if one can redeploy skills and the other cannot.
\end{itemize}
\end{frame}

\begin{frame}{Early choices}
\begin{itemize}
    \item Majors, fields, vocational tracks, and first occupations shape future task exposure.
    \item The main empirical challenge is selection into fields and training paths.
    \item Cutoff and centralized-admissions designs compare students near assignment thresholds.
    \item These designs are strongest for access to early skill pathways, not for all later-career displacement.
\end{itemize}
\end{frame}

\begin{frame}{Training and retraining}
\begin{itemize}
    \item Training can be formal, informal, employer-provided, public, worker-financed, or embedded in job redesign.
    \item Take-up is selected: time, information, motivation, liquidity, and employer support all matter.
    \item A useful design measures take-up, completion, earnings, employment, job quality, and persistence.
    \item Training can reduce inequality if it reaches exposed workers; it can widen inequality if advantaged workers capture it.
\end{itemize}
\end{frame}

\begin{frame}{Within-firm adaptation}
\begin{itemize}
    \item Technology is often implemented through managers, software systems, workflows, and internal labor markets.
    \item Matched employer-worker data can track incumbents before and after adoption.
    \item Key outcomes: retention, exit, earnings growth, task reassignment, occupation switching, and hiring composition.
    \item Main threat: adopters and their workers may already be on different trajectories.
\end{itemize}
\end{frame}

\begin{frame}{Methods and design}
\small
\begin{tabular}{p{0.30\linewidth} p{0.42\linewidth} p{0.16\linewidth}}
\toprule
Design & Identifies well & Main caution \\
\midrule
Major/track cutoff RD & returns to skill pathways & local to cutoff \\
Training evaluation & causal effect of offer or access & take-up selection \\
Matched worker-firm event study & adaptation around adoption & adopter selection \\
Local exposure + admin data & equilibrium training or retirement response & place-level incidence \\
Worker patent/task exposure & worker-level technology incidence & exposure mapping \\
Firm AI/task panels & adoption, training, and job redesign & short horizons \\
\bottomrule
\end{tabular}
\end{frame}

\begin{frame}{What a strong design states}
\begin{enumerate}
    \item Technology shock: robot, AI, software, patent, digital tool, or task automation.
    \item Exposed workers: occupation, firm, industry, place, field, task bundle, or skill history.
    \item Adjustment margin: training, switching, reassignment, migration, retirement, or nonemployment.
    \item Outcome horizon: immediate, medium-run, or career.
    \item Counterfactual and welfare object: earnings, employment, job quality, risk, or broader welfare.
\end{enumerate}
\end{frame}

\begin{frame}{Worker frictions}
\small
\begin{tabular}{p{0.25\linewidth} p{0.42\linewidth} p{0.20\linewidth}}
\toprule
Friction & Adjustment problem & Inequality channel \\
\midrule
Distrust / skepticism & avoid new tools or training providers & past experience differs \\
Inertia / present bias & delay costly learning with future returns & delays are costly \\
Information frictions & do not know which skills pay & networks matter \\
Liquidity and time costs & cannot finance training or search & wealth matters \\
Switching frictions & credentials, housing, family, identity & mobility differs \\
\bottomrule
\end{tabular}
\vspace{0.3cm}
\begin{itemize}
    \item Frictions affect both take-up and the distribution of technology gains.
\end{itemize}
\end{frame}

\begin{frame}{Dynamic worker losses}
\[
Loss_i(H)
=
\sum_{h=0}^{H} \beta^h
\left[
Y_{i,t+h}^{0} - Y_{i,t+h}^{1}
\right]
\]
\begin{itemize}
    \item $Y$ can be earnings, employment, hours, amenities, or job quality.
    \item Fast reemployment can still leave a long earnings loss.
    \item Remaining employed can still mean lower autonomy, lower advancement, or weaker bargaining.
    \item The horizon matters because adjustment costs are dynamic.
\end{itemize}
\end{frame}

\begin{frame}{Week 3 lab}
\begin{itemize}
    \item Reproduce: construct synthetic worker-level technology exposure from task shares.
    \item Diagnose: classify adjustment as staying, training, switching, unemployment, or exit.
    \item Transfer: compare training and retirement responses under local exposure.
    \item Challenge transfer: compare incumbent training and new-hire training under AI adoption.
    \item Goal: see why aggregate skill upgrading can hide incumbent worker losses.
\end{itemize}
\end{frame}

\begin{frame}{Bridge to Week 4}
\begin{itemize}
    \item Week 3 treats workers as the unit that must adjust.
    \item Week 4 moves inside firms.
    \item The next question is how adoption, hiring, management, internal labor markets, and job redesign shape worker adjustment.
\end{itemize}
\end{frame}

\begin{frame}{Takeaways}
\begin{itemize}
    \item Technology changes the return to skills and the risk attached to careers.
    \item Early choices, within-firm retraining, mobility, and exit are separate margins.
    \item Strong empirical designs measure dynamic worker outcomes and confront selection.
    \item Frictions can turn a productivity-enhancing technology into unequal worker incidence.
\end{itemize}
\end{frame}

\end{document}
